\section{Removing the shared variables and completing the trials}\label{cc:section}

The transition in Lemma~\ref{ct:transition} keeps the shared variables
and preserves the conditional records. We now use it finitely many times,
then replace the nonroot shared variables by independent site sampling.
The order matters: the replacement estimate is for one packet, so we first
reduce the trial to a one-packet horizon. The final step approximates the
resulting conditional sampling laws by finite rational tables.

Panchenko's theorem for finite replica symmetry breaking~\cite[arXiv version, Theorem~2]{Panchenko2015FRSB}
removes common nonroot coordinates while retaining common root data, under
spin-tested Ghirlanda--Guerra identities, cavity equations for a modified
Hamiltonian, and an overlap law with finitely many exact values.
Here the overlap cuts partition a possibly continuous distribution.
The recorded conditional laws and the quantitative estimates below provide
the approximation needed to construct finite rational trials.

Panchenko's spin-distribution formula optimizes over a wider invariant
distributional class in its stated even-arity setting; its shared-coordinate
obstruction is explicit after Lemma~6
\cite[arXiv version, Theorem~2 and the Remark after Lemma~6]{Panchenko2013SD}.

\subsection{From recorded tests to probability kernels}

Write $W_1$ for the Wasserstein metric associated with the metric on a
compact type space.  Every iterated type space $K^{(j)}$ in
the structural record~\eqref{ctr:record} is compact and has diameter
at most one. Write $z=(z_0,\ldots,z_d)$ for a complete global path.

\begin{lemma}[Finite tests for the site law]\label{cc:site-law}
Given $e_0>0$ and a finite nonempty list of type spaces, one can choose finite
nonempty collections $\Phi_j$ of $[0,1]$-valued $1$-Lipschitz functions
such that, for probability measures $\nu,S$ on $K^{(j)}$,
\begin{equation}\label{cc:net-bound}
 W_1(\nu,S)
 \le\max_{\phi\in\Phi_j}|\nu(\phi)-S(\phi)|+e_0/2.
\end{equation}
If the structural record holds with
\begin{equation}\label{cc:delta-choice}
 \bigl(\max_j|\Phi_j|\bigr)\delta\le e_0/2,
\end{equation}
then, for each processed block $(b,c]$ and almost every fixed global
path,
\begin{equation}\label{cc:site-law-distance}
 \int W_1\bigl(\nu_{z,u_{0:b}},S_b(z_{0:b},u_{0:b})\bigr)
               \,\dd u_{0:b}\le e_0.
\end{equation}
Here $\nu_{z,u_{0:b}}$ is the conditional law of $S_c$ obtained by
integrating only the fresh site coordinates $u_{b+1:c}$, with all the
global coordinates fixed.
\end{lemma}

\begin{proof}
Kantorovich--Rubinstein duality expresses $W_1$ as the supremum of
integral differences over $1$-Lipschitz functions.  Constants cancel from
such differences; since the space has diameter at most one, each test
can be shifted into $[0,1]$.  The class of these bounded Lipschitz
functions is compact in the uniform norm.  A finite $e_0/4$-net gives
Equation~\eqref{cc:net-bound}, since replacing a test by a net point
changes the difference of its two integrals by at most $e_0/2$.

For a fixed test, the structural record and Cauchy--Schwarz give
\[
 \int |\nu_{z,u_{0:b}}(\phi)-S_b(z_{0:b},u_{0:b})(\phi)|
                 \,\dd u_{0:b}\le\delta.
\]
Bound the maximum over tests by their sum, integrate
Equation~\eqref{cc:net-bound}, and use
Equation~\eqref{cc:delta-choice}.  This proves
Equation~\eqref{cc:site-law-distance} with its almost-everywhere global
quantifier intact.
\end{proof}

\subsection{Finite stopping}

\begin{proposition}[A terminal one-packet trial]\label{cc:terminal}
Fix $U>P(a,\beta)$, $0<\eta,m_0<1$, and an integer $J$ such that
\begin{equation}\label{cc:stopping-depth}
 (1-\eta)^J<m_0.
\end{equation}
Choose in advance finite nonempty test sets on $K^{(j)}$, $0\le j\le J$,
and a structural tolerance $\delta>0$.  For every $\rho>0$ there is a
trial with
\[
 G_1\le U+\rho,
\]
at most $J$ processed blocks, and depth at most $2J+1$, which satisfies
the chosen structural records and has one of the following forms:
\begin{enumerate}
 \item the processed blocks start at the root boundary $0$;
 \item the processed blocks start at boundary $1$, and the remaining
       first exponent satisfies $0<m<m_0$.
\end{enumerate}
Every processed block has exponent width at most $\eta$.
\end{proposition}

\begin{proof}
Lemma~\ref{ctr:reservoir} supplies trials for arbitrarily large
integer horizons $N$, initially of depth one and first exponent $1$,
with $G_N/N\le U$.  There are initially no processed blocks.  A trial
is terminal if its blocks start at the root or if its first unprocessed
exponent is less than $m_0$.

Suppose a sequence at a given stage contains no terminal trial.  All its
first exponents are then at least $m_0$.  Depth has increased by at most
two at each preceding transition, so there are only finitely many
possible patterns of block boundaries.  Pass to an unbounded horizon
subsequence with a common pattern, and then to a subsequence on which
the exponents converge.  Lemma~\ref{ct:transition}, with its
fixed tests and tolerance, produces for each fixed positive integer $k$
an output trial satisfying $G_k/k\le U$.  It preserves the old records
and adds one processed block.  Either the output is terminal, or its
new first exponent is at most $(1-\eta)$ times the limiting old first
exponent.

This use of the transition involves nested limits, not a uniform limit
in $k$.  First fix $k$; then take the input horizon to infinity as
required by the transition; finally retain an output for that $k$.
Doing this separately for $k=1,2,\ldots$ gives the next unbounded
horizon sequence if no terminal output occurs.  After $j$ transitions
every still-active output has first exponent at most $(1-\eta)^j$.
Equation~\eqref{cc:stopping-depth} therefore forces termination after
at most $J$ transitions.  The resulting terminal trial has some finite
horizon, at most $J$ blocks, and depth at most $2J+1$.

Reduce its horizon to one using the packet telescope of
Lemma~\ref{ctr:posterior}, without a Gaussian perturbation.  There is
no remainder when the blocks have size one.  The sum of the averaged
one-packet posterior differences is the original $G_k$, so one step,
and then base root data outside the conditional null sets, have
posterior value at most $U+\rho$.  This selection is made before fresh
packet data are sampled.  The conditional laws generated by the
posterior densities remain equivalent to the old laws, and the old
site-law functions are kept rather than recomputed.  Hence all the
almost-everywhere records persist.  Neither the exponents nor the
terminal boundary changes.  This gives the asserted one-packet trial.
\end{proof}

\subsection{Removing shared variables after stopping}

\begin{lemma}[Collapse by the recorded kernels]\label{cc:kernel-collapse}
Consider a terminal trial from Proposition~\ref{cc:terminal}, whose
records imply Equation~\eqref{cc:site-law-distance}.  Fix a count cutoff
$D_0\ge1$ and put
\[
 C=\log2+\beta D_0,\qquad M=3D_0,
 \qquad
 R=JM e^\beta e^{2C(2J+2)}e_0.
\]
There is a site-only trial, with weakly increasing exponents in $(0,1]$,
such that
\begin{equation}\label{cc:collapse-error}
 |G_1^{\mathrm{site}}-G_1|
 \le 2K(C)\eta+2R+C^2e^{2C}m_0+2\tau_A+2\tau_B,
\end{equation}
where
\[
 \tau_A=\E[(\log2+\beta D_A)\1_{\{D_A>D_0\}}],\qquad
 \tau_B=\E[(\log2+\beta D_B)\1_{\{D_B>D_0\}}],
\]
with $D_A\sim\Poi(3a)$ and $D_B\sim\Poi(2a)$.
In the root-terminal case the term containing $m_0$ is unnecessary.
The new message kernels are independent of the fresh packet counts
and signs.
\end{lemma}

\begin{proof}
We first condition on either endpoint's fresh count and sign data with
$D\le D_0$.  Its loss $X$ is bounded in absolute value by $C$, and it
uses at most $M$ distinct reservoir occurrences.  Simultaneously replace
all exponents in each processed block $(b,c]$ by $t_c$.  The resulting
vector is still nondecreasing and is within $\eta$ of the original.
Lemma~\ref{cs:sensitivity} bounds the change by $K(C)\eta$.
A block with constant exponent is one conditional $T_{t_c}$ operation
over all its coordinates, by the tower property.

Define deterministic functions of vectors of single-site types,
one coordinate for each reservoir occurrence, working upwards through
the block boundaries.  At the last boundary put
$F_d(\mathbf S_d)=X$.  For a block $(b,c]$ put
\begin{equation}\label{cc:kernel-recursion}
 F_b(\mathbf s)=\frac1{t_c}\log\int
       \exp\bigl(t_cF_c(\mathbf u)\bigr)\prod_i s_i(\dd u_i).
\end{equation}
These functions depend on the fixed count and sign data, but the
probability kernels $S_b$ do not.  The functions $F_b$ are bounded in
absolute value by $C$.

Use the sum metric on each finite product of type spaces.  The terminal
loss is $e^\beta$-Lipschitz: a clause factor is bounded below by
$1-w=e^{-\beta}$, the product of at most three $[0,1]$ messages changes
by at most the sum of their changes, and the log-sum-exp for the two
branches of an A packet preserves the largest branchwise log error.
Thus we may take $L_d=e^\beta$ as a Lipschitz constant for $F_d$.
Coordinatewise couplings in Equation~\eqref{cc:kernel-recursion} give
\begin{equation}\label{cc:lipschitz-recursion}
 L_b\le e^{2C}L_c.
\end{equation}
Indeed $v\mapsto e^{t_cv}$ has Lipschitz constant at most $t_ce^C$ on
$[-C,C]$, and the divided logarithm on the possible integral values
has Lipschitz constant at most $e^C/t_c$.  The same two bounds show that
a mean absolute error in an inside value propagates through a block
by at most the factor $e^{2C}$.

Let $Y_b$ denote the actual effective value at boundary $b$ after
flattening the processed blocks' exponents.  Thus $Y_d=X$, and across
a block $(b,c]$ the value $Y_b$ is the conditional $T_{t_c}$ average
of $Y_c$ over the block's global and site coordinates.  Define
\[
 E_b=\E_{\mathrm{base}}
       \bigl|Y_b-F_b(\mathbf S_b)\bigr|.
\]
Here the expectation uses the current trial law on the global history
and all occurrences' fresh site histories through $b$, before any
exponential reweighting by this fresh packet.  This law already includes
any old-packet posterior used to obtain the trial; the packet count and
signs remain fixed.  In particular $E_d=0$.  We claim that
\begin{equation}
 E_b\le e^{2C}E_c+e^{2C}L_cMe_0.
 \label{cc:one-block-error}
\end{equation}

First replace $Y_c$ inside the block average by
$F_c(\mathbf S_c)$.  The mean-error propagation bound above gives
the first term of Equation~\eqref{cc:one-block-error}.  To bound the remaining error, fix the
ancestor histories through $b$ and the continuation globals through
$c$.  Under this conditioning, the occurrences' fresh site
continuations are independent; write $\nu_i$ for the resulting law
of $S_c^i$.  A product of Wasserstein couplings and the Lipschitz
constant of $F_c$ give
\[
 \left|\int e^{t_cF_c(\mathbf u)}\prod_i\nu_i(\dd u_i)
       -\int e^{t_cF_c(\mathbf u)}\prod_iS_b^i(\dd u_i)\right|
 \le t_ce^C L_c\sum_i W_1(\nu_i,S_b^i).
\]
The second integral depends only on the ancestor histories.  Average
this inequality over the continuation globals, apply the
$e^C/t_c$ Lipschitz bound for the divided logarithm, and then average
over the ancestor histories.  For almost every fixed global path,
Equation~\eqref{cc:site-law-distance} bounds the integral over each
occurrence's site ancestors by $e_0$.  The continuation-global law
is independent of these fresh site generators, including when the
trial was obtained by a posterior change of law.  Fubini's theorem
therefore bounds the final average by $e^{2C}L_cMe_0$ and proves Equation~\eqref{cc:one-block-error}.

There are at most $J$ processed blocks, and
$L_c\le e^{2CJ}e^\beta$ at every boundary.  Iterating Equation~\eqref{cc:one-block-error} from
$E_d=0$ bounds their accumulated error.  If the starting boundary
is $1$, propagating once through the remaining first-level transform
adds at most one factor $e^{2C}$.  In either case the error is at most
\[
 JM e^\beta e^{2C(2J+2)}e_0=R.
\]

If the processed blocks begin at the root, the resulting recursion
already describes a site-only trial. Retain the common root $z_0$
and sample each occurrence's initial type $S_0(z_0,u_0^i)$ using its
own fresh root-site uniform. At one level of exponent $t_c$ for each
old block, draw each site's next type independently from its preceding
type, and continue to the message pair.  Probability kernels on these compact metric spaces
can be sampled by measurable functions of fresh uniform variables.
This construction evaluates Equation~\eqref{cc:kernel-recursion}.

If the blocks begin at boundary $1$, there is still the operation
$T_mF_1(\mathbf S_1)$, where $0<m<m_0$.  Conditional on the root data,
factor this operation as
\[
 T_m^{z_1}T_m^{\mathrm{sites}}F_1(\mathbf S_1).
\]
The inside quantity is in $[-C,C]$.  By
Equation~\eqref{cs:small-exponent}, replacing the outside
$T_m^{z_1}$ by ordinary expectation changes its value by at most
$\tfrac12 C^2e^{2C}m_0$.  Move $z_1$, with this ordinary sampling law,
into the root global data.  Retain the level-$1$ site transform and
the independent block draws just constructed.  This is again
site-only.  Its kernels were constructed from the recorded types
alone, not from any fresh count or sign data.

For counts greater than $D_0$, both the old and new endpoint transforms
are bounded in absolute value by $\log2+\beta D$, regardless of their
kernels or exponents.  Their contribution to the difference is thus
at most twice the corresponding tail expectation.  Summing the
bounded-count estimates for the two endpoints proves
Equation~\eqref{cc:collapse-error}.
\end{proof}

\subsection{Rational finite descriptions}

\begin{theorem}[Finite rational completeness]\label{cc:complete}
\label{cc:rational}
For every $a>0$ and finite $\beta>0$,
\begin{equation}\label{cc:finite-infimum}
 P(a,\beta)=\inf_{\mathcal Q}G_1,
\end{equation}
where $\mathcal Q$ consists of site-only trials of finite depth whose
exponents are strictly increasing rationals in $(0,1)$ and whose
message pair is a rational-valued step function on a finite rational
rectangular grid in its root and site uniform variables.  In
particular, this is a countable class with an effective enumeration.
\end{theorem}

\begin{proof}
The upper bound of Proposition~\ref{ctr:upper} gives
$P(a,\beta)\le G_1$ for every member of $\mathcal Q$.  We prove the
reverse inequality by giving the order of all approximation choices.

Fix $\epsilon>0$ and put $U=P(a,\beta)+\epsilon$.  First choose a
count cutoff $D_0\ge1$ with $\tau_A,\tau_B\le\epsilon$, and set
$C=\log2+\beta D_0$.  Choose $\eta,m_0\in(0,1)$ so that
\[
 K(C)\eta\le\epsilon,\qquad
 \tfrac12 C^2e^{2C}m_0\le\epsilon.
\]
Next choose $J$ satisfying Equation~\eqref{cc:stopping-depth}, and
only then choose $e_0>0$ so small that
\[
 J(3D_0)e^\beta e^{2C(2J+2)}e_0\le\epsilon.
\]
Choose the finite test sets on all type spaces through $K^{(J)}$ by
Lemma~\ref{cc:site-law}, and finally choose $\delta$ satisfying
Equation~\eqref{cc:delta-choice}.  In particular, neither the test
precision nor the allowed number of transitions is chosen in response
to the eventual hierarchy depth.  The overlap widths used inside each
transition are chosen only after these test sets and $\delta$, as in
Lemma~\ref{ct:transition}.

Proposition~\ref{cc:terminal}, with $\rho=\epsilon$, gives a
terminal one-packet trial of value at most $P+2\epsilon$.
Lemma~\ref{cc:kernel-collapse} turns it into a site-only trial whose
value is at most $P+12\epsilon$: the exponent changes contribute at
most $2\epsilon$, the kernel replacements at most $2\epsilon$, the
possible small-exponent replacement at most $2\epsilon$, and the
two tails at most $4\epsilon$.

It remains to obtain a finite rational description.  Approximate the
finite weakly increasing exponent vector by strictly increasing
rational vectors in $(0,1)$.  This is possible even if some original
exponents coincide or equal $1$.  For fixed count data,
Lemma~\ref{cs:sensitivity} gives convergence of the
transform values.  The packet bound of
Equation~\eqref{ctr:packet-bound}, uniform over the approximations,
then gives convergence after averaging the Poisson counts and signs.

Encode the root randomness by a single uniform coordinate.  After the
independent conditional type samplings have likewise been encoded by
uniform coordinates, the message pair is a measurable function of
finitely many independent uniforms: the root global, the root site
coordinate, and the finitely many site-level coordinates.  Approximate
this pair in $L^1$ by rational-valued step functions on finite dyadic
rectangular grids.  For example, conditional expectations on the
successive dyadic grids converge in $L^1$, and subsequent rational
quantization preserves the range $[0,1]^2$ and the convergence.
The shared root coordinate is retained as shared; no independence
across sites is introduced there.

For fixed counts and signs, the leaf Lipschitz estimate in the proof
of Lemma~\ref{cc:kernel-collapse} bounds the mean leaf error by
$e^\beta$ times the sum of the $L^1$ message errors over the finitely
many occurrences.  The mean-error propagation bound through the
fixed finite hierarchy proves convergence of both endpoint values.
Once more, the root linear-in-count bound makes the Poisson tails
uniform, so their fully averaged $G_1$ values converge.  Thus a member
of $\mathcal Q$ has value at most $P+13\epsilon$.

Since $\epsilon$ is arbitrary, Equation~\eqref{cc:finite-infimum}
follows.  Finite depths, finite grids, rational exponents, and finite
rational tables all have finite encodings and can be enumerated by a
single ordinary algorithm.  These are the trials used in the
certificate search; none of the compactness limits or tolerance
choices in the existence proof is input to that search.  The larger class of strict measurable site-only trials contains
$\mathcal Q$ and obeys the same upper bound, so its infimum also equals
$P(a,\beta)$.
\end{proof}
